% --- TCOLORBOX ---
\usepackage[many]{tcolorbox}

\tcbset{
	stilebase/.style={
		enhanced,
		breakable, 
		colback=white,
		coltitle=white,
		toptitle=1.2mm,
		bottomtitle=1.2mm,
		boxrule=1pt, 
		separator sign={:},
		before skip=12pt, 
		after skip=12pt,
		fonttitle=\bfseries\large\raggedright
	}
}

% \label nel corpo di un box rimanda al numero del box (non a quello della sezione)
\makeatletter
\newcommand{\tcbsetcurrentlabel}{\protected@edef\@currentlabel{\thetcbcounter}}
\makeatother

% Fabbrica di box per OCP (Open/Closed Principle)
\newcommand{\newmanualbox}[3]{%
    % #1 = environment name
    % #2 = title prefix
    % #3 = additional tcb keys
    \newtcolorbox[auto counter,number within=chapter,list inside=#1]{#1}[1][]{%
        stilebase,
        adjusted title={#2 \thetcbcounter\if\relax\detokenize{##1}\relax\else: ##1\fi},
        before upper={\tcbsetcurrentlabel},
        #3
    }%
}

% Ambiente per definizioni
\newmanualbox{definizione}{Definizione}{colframe=black!80!white}

% Ambiente per esempi
\newmanualbox{esempio}{Esempio}{colframe=black!60!white, boxrule=0.6pt, before upper app={\interlinepenalty=10000\relax}}

% Ambiente per regole
\newmanualbox{regola}{Regola}{colframe=black!80!white}

% Ambiente per errori comuni
\newmanualbox{errore}{Errore}{colframe=black!95!white}

% Ambiente per fisica
\newmanualbox{fisica}{Con la fisica}{colback=black!8!white,colframe=black!50!white,boxrule=0.4pt,fontupper=\small,sharp corners}

% Ambiente per matematica
\newmanualbox{matematica}{Con la matematica}{colback=black!8!white,colframe=black!50!white,boxrule=0.4pt,fontupper=\small,sharp corners}

% Ambiente per osservazioni
\newmanualbox{osservazione}{Osservazione}{colback=black!2!white,colframe=black!40!white,boxrule=0.4pt,fontupper=\small\itshape,sharp corners}
